% V2 PDF物理页73--74 原书思考题14.1--14.11；图14.13等待专家校核。
\section*{第14章\quad 静电场中的导体：思考题解析}
\paragraph{思考题 14.1}
任意形状的理想导体，只要达到\emph{静电平衡}，内部自由电荷停止净漂移，必须 \(\boxed{\boldsymbol E_{\rm metal}=0}\)，否则载流子持续运动。此结论不要求球形；尖端附近外部场可能很强，但不改变导体体内零电场。

\paragraph{思考题 14.2}
点电荷在腔中心时，内表面感应 \(-q\)，外表面总电荷 \(Q+q\)。球壳为等势体，壳外电势（以无穷远零） \(\varphi_{\rm out}(r)=(Q+q)/(4\pi\epsilon_0r)\)；腔内 \(r<R\) 有点电荷场加上边界决定的常数。因此不能把已带电孤立壳的\(\varphi_0\)在整个空间\emph{不变}地加上点电荷势；插入q改变了壳面电荷分布。

\paragraph{思考题 14.3}
新增电子数 \(N\) 使金属球的静质量至少增加 \(Nm_e\)（如果电子最终仍在球上）。但系统总能量变化还包括电场储能、电子束缚能和放出的热/辐射；若严格使用 \(E=mc^2\) 衡量整个封闭系统质量，不能只加 \(Nm_e\)。在通常宏观力学精度下，球质量\emph{增加}而非因为负电而减少。

\paragraph{思考题 14.4}
点电荷居中心时球对称强制内表面 \(\sigma_{\rm inner}=-q/(4\pi R_{\rm in}^2)\) 均匀；孤立球壳外表面若无外场也均匀，\(\sigma_{\rm outer}=(Q+q)/(4\pi R_{\rm out}^2)\)。移离中心破坏内腔球对称，内表面面电荷密度不均匀，但总量仍\(\boxed{Q_{\rm inner}=-q}\)。外壳若球形、无外场仍能保持球对称外电势。

\paragraph{思考题 14.5}
封闭导体壳处于静电平衡时金属体内 \(\boldsymbol E=0\)，外部电荷引起的表面自由电荷重排，使封闭腔内满足同一边界等势条件。若腔内无电荷，则由唯一性定理腔内 \(\boldsymbol E=0\)。若腔内有固定电荷，其电场与导体壁边界的解不因外部静电源改变；腔的绝对电势常数可随孤立导体电势改变。

\paragraph{思考题 14.6}
取横跨导体界面的薄 Gauss 小柱，内部 \(E_n=0\)、横向面积为A，则 \(\epsilon_0 E_{\perp,\mathrm{out}}A=\sigma A\)，故 \(\boxed{E_{\perp,\mathrm{out}}=\sigma/\epsilon_0}\)。邻近外电荷会重新分配局部\(\sigma\)与外侧总场，但在静电理想导体表面、无界面额外表面电荷和真空外界条件下，此\emph{局部边界关系仍成立}。

\paragraph{思考题 14.7}
若原题断言对\emph{任何}两带电导体“至少有一处面电荷同号”，需要补充对总电荷符号、外部电荷和边界的假设。对两个都具有正净电荷的孤立导体，若某一个表面处处非正，则面积积分其净电荷不可能为正，因此它必有 \(\sigma>0\) 的区域；若两者净电荷异号或另有外场，则不能得到不依赖条件的共同符号结论。导体整体 \(\int\sigma dA=Q\) 只是积分约束，不决定每个面元的符号。

\paragraph{思考题 14.8}
孤立无限电荷薄面场由对称性分配到两侧：每侧 \(\boxed{|E|=\sigma/(2\epsilon_0)}\)。导体表面则一侧为金属内部 \(E=0\)，另一侧承担全部边界法向场跳变，故 \(\boxed{|E_{\rm out}|=|\sigma|/\epsilon_0}\)。两种几何的表面电荷分别承担不同的总边界条件，不能把薄绝缘板模型当导体模型。

\paragraph{思考题 14.9}
两个无限大平行导体板若外部电场消失，Gauss边界关系与平面对称性推出面对内侧的电荷为 \(\sigma\) 与 \(-\sigma\)。有限板有边缘电场、外表面可存电荷，仍须各自满足静电等势与总电荷守恒。一般只能说在大平板的\emph{远离边缘}区域，内侧电荷密度近似等量异号，不能无条件要求\emph{所有}表面点配对。

\paragraph{思考题 14.10}
孤立导体球的电势由其自身总电荷与外部电荷的相互感应共同决定；点电荷q移到壳外改变电荷分布，一般使球电势发生变化，不能仅靠q不变就判定电势不变。若球接地则电势由大地固定为0，外部点电荷使\emph{导体上总感应电荷}改变。若球空腔中的点电荷移出且孤立球原净电荷Q固定，则导体总电荷Q不会随移动改变。

\paragraph{思考题 14.11}
最初C、D异号电荷靠近导致A/B相连导体极化。把C、D导通后，它们之间电荷重新分配，外部电势边界与场发生变化，从而A/B上的感应面电荷也\emph{一般重新分布}。若C+D总电荷为零且二者构型允许完全中和，外场可能显著减弱，A/B感应电荷趋零；但如果C/D还受其他电荷影响或非对称分布，不能仅凭“异号相吸”声称A/B上的所有感应电荷必然不变或全消失。
